% =========================
% Cas pratique 8
% =========================

\section[Cas 8]{Cas pratique 8 : Interdiction provisoire}

%--------------------------
\begin{frame}{Question 1}

\begin{block}{}
L'art. \textbf{L615-3} CPI dispose que "peut saisir en référé
la juridication civile compétente" ou "la juridication civile
compétente peut également ordonner [...] sur requête".
\end{block}

\begin{itemize}
  \item En l'espèce, Phyto souhaite l'une quelconque des
  interdictions provisoire.
  \item En conclusion, Phyto a la possibilité de s'adresse
  au juge des référés ou au juge de mise en l'état.   
\end{itemize}

\end{frame}
%--------------------------


%--------------------------
\begin{frame}{Question 2}

\begin{block}{}
L'art. \textbf{L615-3} al. 1 3e phrase CPI dispose que
"preuve raisonnablement accessible au demandeur et
vraisemblable qu'il est porté atteinte à ses droits"
\end{block}

\begin{itemize}
  \item En l'espèce, Phyto estime qu'il est porté
  atteinte à certaines revendications de son brevet.
  On suppose qu'une preuve raisonnablement accessible
  au demandeur est disponible. Il n'y a pas d'urgence
  apparente.
  \item En conclusion, la procédure en référé est
  préférée.
\end{itemize}

\end{frame}
%--------------------------

%--------------------------
\begin{frame}{Question 3}

\begin{block}{}
L'art. \textbf{L615-3} al. 1 3e phrase CPI ne dispose pas
d'une condition sur le statut du brevet.
\end{block}

\begin{itemize}
  \item En l'espèce, le brevet est l'objet d'un recours.
  \item En conclusion, une mesure provisoire n'est pas
  interdite.
\end{itemize}

\end{frame}
%--------------------------

%--------------------------
\begin{frame}{Question 4}

\begin{block}{}
\textbf{Principe de proportionnalité} (Ordonnance JME TGI Paris, 19 avril 2013, page 5) :
"raisonnablement pas ignorer la nullité [des revendications]". 
\end{block}

\begin{itemize}
  \item En l'espèce, le brevet est maintenu après la 1ère instance sous forme limitée.
  La validité du brevet est vraisemblable.
  \item En conclusion, les chances de succès d'une demande d'interdiction provisoire sont bonnes.
\end{itemize}

\end{frame}
%--------------------------

%--------------------------
\begin{frame}{Question 5}

\begin{block}{}
Art. 490 CPC (référé) ou Art. 496 CPC (requête) + A539 CPC (principe)
\\ Exception (Ordonnance TGI Paris, 15 avril 2016): "indépendamment de
l'absence de caractère suspensif de ce recours".
\end{block}

\begin{itemize}
  \item Par hypothèse, une interdiction provisoire est ordonnée.
  Hypothèse 1: en référé; Hypothèse 2: sur requête.
  \item En conclusion, dans les 2 cas, un recours est possible,
  mais pas suspensif.
\end{itemize}

\end{frame}
%--------------------------